\section{La perspectiva de la producción en serie: tecnología, calidad, costo y responsabilidad deben evaluarse en conjunto}

Antes se trataron la ruta tecnológica y la simplificación de la organización; esta sección agrega la perspectiva de la producción en serie. Liu Xianming (刘先明) subrayó varias veces que prefiere encontrar en los problemas reales lo que debe hacerse en la siguiente etapa. El responsable de conducción autónoma no es solo un responsable de investigación: también debe atender el producto, el negocio, la calidad, el hardware, el costo y la seguridad de los usuarios. Esta perspectiva explica por qué un fabricante de automóviles no puede copiar sin más el ritmo de un equipo de investigación.

\subsection{La producción en serie no es la etapa final de la investigación, sino una fuente de problemas}

Los equipos de investigación suelen definir primero la tarea y después buscar los datos y el modelo; el equipo de producción en serie enfrenta los problemas que los usuarios reales encuentran todos los días. Cada toma de control por parte del conductor, cada maniobra incómoda, cada error de juicio y cada queja de un usuario en la carretera se convierte en un problema que la siguiente iteración del modelo y de la ingeniería debe resolver. La producción en serie no es la fase de despliegue que sigue al término de la investigación, sino la entrada por la que surgen continuamente nuevos problemas de investigación.

\begin{knowledgebox}{Cómo los problemas de la producción en serie retroalimentan al modelo}
\begin{tabularx}{\textwidth}{p{0.22\textwidth}p{0.34\textwidth}X}
\toprule
Señal de la producción en serie & Qué hacer antes de que llegue al modelo & Valor de la retroalimentación \\
\midrule
Toma de control & Localizar la causa de la toma de control, extraer el fragmento y etiquetar el contexto & Encontrar los puntos ciegos del modelo y el límite de las reglas. \\
Incomodidad del usuario & Relacionarla con la maniobra, la velocidad, la distancia entre vehículos y el entorno & Mejorar la comodidad y el estilo de conducción. \\
Accidente/riesgo & Atribuir la responsabilidad y revisar la cadena de percepción/predicción/planeación & Formar muestras de seguridad y conjuntos de evaluación. \\
Restricciones de hardware & Registrar los límites de capacidad de cómputo, sensores, latencia y costo & Orientar la destilación, la cuantización y la poda. \\
\bottomrule
\end{tabularx}
\end{knowledgebox}

\subsection{Por qué «lo simple es bello» no significa eliminar sin criterio}

Simplificar no es eliminar todas las protecciones, sino quitar las estructuras repetidas, obsoletas o que obstaculizan el scaling. Un fabricante de automóviles todavía debe conservar la redundancia de seguridad, la validación de calidad y el proceso de liberación a producción. La simplicidad que subraya Liu Xianming consiste en que el proceso de desarrollo sea más corto, que haya menos objetivos y que los problemas salgan a la luz más rápido, no en sacrificar los límites de seguridad.

\begin{warningbox}{Simplificar no equivale a precipitarse}
Si, por perseguir el enfoque de extremo a extremo, se descuidan la seguridad, las pruebas y la calidad de la producción en serie, la simplificación se convierte en precipitación. La simplificación verdadera reduce la complejidad inútil y, al mismo tiempo, conserva las restricciones de seguridad necesarias.
\end{warningbox}

\subsection{Resumen de la sección}

La ruta de producción en serie de Physical AI exige que la tecnología y la responsabilidad avancen juntas. Los modelos deben ser más grandes y los datos más abundantes, pero cada liberación a producción debe superar la evaluación de seguridad, calidad, costo y experiencia de usuario.
